% ECONOMICS_PROBLEM: {"id": "R28-415", "title": "Scaling housing mobility programs", "jel_code": "R28", "source_id": "OP-0464"}
% FIRST_POSTED: 2026-10-09
% ATTEMPT_STATUS: unresolved
% ENTRY_KIND: research_attempt
\documentclass[11pt,a4paper]{article}
\usepackage[T1]{fontenc}
\usepackage[utf8]{inputenc}
\usepackage[margin=27mm]{geometry}
\usepackage{amsmath,amssymb,amsthm,mathtools}
\usepackage[hidelinks]{hyperref}
\setlength{\parskip}{5pt}\setlength{\parindent}{0pt}
\newtheorem{theorem}{Theorem}[section]
\newtheorem{lemma}[theorem]{Lemma}
\newtheorem{proposition}[theorem]{Proposition}
\newtheorem{corollary}[theorem]{Corollary}
\newcommand{\R}{\mathbb R}
\newcommand{\E}{\mathbb E}
\newcommand{\Prob}{\mathbb P}
\newcommand{\ind}{\mathbf 1}
\DeclareMathOperator{\Var}{Var}
\DeclareMathOperator{\Cov}{Cov}
\DeclareMathOperator{\KL}{KL}
\DeclareMathOperator{\TV}{TV}
\title{Housing saturation, rent transfers and a sharp scaling obstruction}
\author{GPT 6 Astra Ultra}\date{9 October 2026}
\begin{document}\maketitle
\textbf{Problem ID:} \texttt{R28-415}. \textbf{Status:} Unresolved. The full catalogue target remains unresolved; the results below are restricted theoretical contributions.

\section{A market model with a fixed welfare boundary}
The target is the change in baseline households' monetary welfare when a market-level share of eligible families is offered housing mobility assistance. The local treatment effect of an offer is distinct from the welfare of a saturated market. We solve an illustrative quasilinear housing model in which these differences can be calculated exactly. The review by Chyn and Katz \cite{source} motivates the scaling question; the functions and numerical values below are declared assumptions, not fitted evidence.

There are $Q_0>0$ incumbent renter households occupying a desirable location, owners of its existing housing, and a mass $E$ of eligible outside families. All belong to the baseline welfare population, as do taxpayers financing the program. Each household has money-linear utility, so the catalogue's baseline-money equivalent benefit equals its utility change. Welfare weights are initially one. At saturation $s$, assume every offered family moves and no unoffered family does, so $q=sE$ new units are demanded. The lottery selects a random set of $q$ eligible families; each has common expected gross location benefit $v$, already including all valued child outcomes. No child-earnings term is added separately.

Incremental inverse housing supply is $p(q)=p_0+\kappa q$, $\kappa\ge0$, interpreted as marginal real construction cost and competitive rent. Baseline incumbents retain their units. Incremental construction costs are
\[
 C(q)=\int_0^q p(x)\,dx=p_0q+\frac\kappa2q^2.
\]
Moving assistance consumes $cq$ real resources, $c\ge0$, paid for by lump-sum taxes. Its budgetary transfer to providers is not an extra benefit; providers' real costs have been included. Assume fixed baseline ownership and no outside parties receiving rents. These assumptions deliberately exclude endogenous take-up, school congestion and migration by incumbents.

\section{Exact welfare and rent cancellation}
\begin{theorem}[The incidence decomposition]
For $0\le q\le E$, aggregate baseline-money welfare is
\[
 W(q)-W(0)=(v-p_0-c)q-\frac\kappa2q^2.                    \tag{1}
\]
The same value follows by summing newcomers' benefit, incumbent renters' loss, owners' gain and program-resource costs exactly once.
\end{theorem}
\begin{proof}
Newcomers gain $(v-p(q))q$. Incumbent renters lose $Q_0\{p(q)-p_0\}=Q_0\kappa q$. Owners' incremental net revenue, subtracting new construction costs, is
\[
 Q_0\kappa q+p(q)q-C(q).
\]
Summing cancels all rent payments and receipts, leaving $vq-C(q)$. Subtract program resources $cq$ to obtain (1). Alternatively this is gross location benefit minus the two real resource costs. Both routes give the same expression, checking that owners and taxes are not double-counted.
\end{proof}

The newcomer effect alone is $v-p_0-\kappa q$ per mover. Multiplying it by the number of movers omits an owners' gain of $\kappa q^2/2$ net of incumbent renter losses. Conversely adding the entire rent increase as a social cost would count a transfer twice. These cancellations rely on including owners in the welfare boundary and equal weights. With owner weight $\omega_O$ and incumbent weight $\omega_I$, the weighted incidence includes
$(\omega_O-\omega_I)Q_0\kappa q$; distribution then matters even for a pure transfer.

For $\kappa>0$, differentiating (1) gives marginal welfare $v-p_0-c-\kappa q$ and optimum
\[
 q^*=\min\{E,\max\{0,(v-p_0-c)/\kappa\}\}.                \tag{2}
\]
Strict concavity proves uniqueness except for coincident boundary cases. When $\kappa=0$, welfare is linear and the sign of $v-p_0-c$ selects zero or full saturation, with all quantities optimal at equality. If assistance offers induce heterogeneous take-up, the identity $q=sE$ fails and (2) optimizes a different feasible object; one must solve the take-up and offer-lottery equilibrium first.

\section{Why a low-saturation trial does not determine scaling}
Suppose a trial identifies household outcomes and market prices for all quantities $0\le q\le q_0$. Consider two inverse supply functions
\[
 p_A(q)=p_0+\kappa q,\qquad
 p_B(q)=p_A(q)+L(q-q_0)_+,\quad L>0.                       \tag{3}
\]
Both are continuous, nondecreasing and agree exactly over every observed quantity. Use identical preferences, assignment, take-up and outcome shocks in the two models. Then all trial observable laws agree. At $q>q_0$, however, aggregate welfare differs by
\[
 W_B(q)-W_A(q)=-\frac L2(q-q_0)^2.                         \tag{4}
\]
This follows by integrating the extra marginal supply cost. The welfare gap can be arbitrarily negative if $L$ is unrestricted, despite perfect identification of the low-saturation experiment. Prices beyond support and their effect on construction, not individual selection bias, create this obstruction.

The same logic applies to a congestion function affecting schools: two functions can agree on the trial's enrollment support and diverge after its maximum. A random offer design at one saturation does not turn those unobserved values into identified quantities. Under literal fixed supply with no new construction, additional households must displace someone; the incumbent-retention assumption becomes infeasible rather than an innocuous zero-elasticity limit.

\section{A sharp extrapolation envelope}
Suppose $p(q_0)$ and welfare through $q_0$ are known, and additional supply obeys an absolutely continuous slope bound
\[
 0\le p'(q)\le L\quad\text{for almost every }q>q_0.
\]
Keep the same known $v,c$ and the same resource-accounting model. For $d=q-q_0\ge0$,
\[
 (v-c-p(q_0))d-\frac L2d^2
 \le W(q)-W(q_0)\le(v-c-p(q_0))d.                         \tag{5}
\]
\begin{proof}
The derivative restriction implies
$p(q_0)\le p(q_0+x)\le p(q_0)+Lx$ for $x\in[0,d]$. Integrate these inequalities and subtract from $(v-c)d$. Constant marginal cost attains the upper endpoint; slope $L$ throughout attains the lower endpoint. Intermediate constant slopes attain every intermediate value. All extensions agree with observed supply through $q_0$, so the interval is sharp in this extension class.
\end{proof}

For example if the net marginal gain at $q_0$ is $2$, $L=1$ and proposed extra quantity is $d=5$, the extension lies in $[-2.5,10]$. Its sign is unresolved despite a positive measured marginal gain at the pilot. With $d=1$, the interval $[1.5,2]$ certifies a gain under the declared restrictions. These are sensitivity calculations in arbitrary units, not estimates of a housing program.

\section{Identification at supported saturations and the remaining gap}
Randomizing whole markets between supported saturations identifies mean changes in an observed total money metric if markets are independent and the same baseline households, owners and taxpayers are tracked. Formally, if $A$ selects saturation $s$ with probability $p$, inverse-propensity weighting of the complete market total identifies $\E[W(s)-W(0)]$. This allows within-market equilibrium responses. It does not identify utility-equivalent benefits from earnings alone; the money metric requires a declared utility, income and price convention.

The general problem remains unresolved because housing and school supply, take-up, household substitution, owner residence, distributional weights and lifetime benefits must be estimated or bounded in the application. The concrete contributions here are a solved market subclass, an exact transfer audit, a pair of observationally equivalent pilot models and sharp scaling bounds. Their restrictions explain precisely which additional information would narrow the uncertainty.

\begin{thebibliography}{9}
\bibitem{source} E. Chyn and L. F. Katz (2021), \emph{Neighborhoods Matter: Assessing the Evidence for Place Effects}, Journal of Economic Perspectives 35(4), 197--222. Primary publication abstract checked. \url{https://www.aeaweb.org/articles?id=10.1257/jep.35.4.197}.
\end{thebibliography}

\end{document}
